\documentclass[11pt]{article}
\usepackage{marcoreport}
\usepackage{listings}
\usepackage{xcolor}

% Code listing style
\lstset{
  language=Python,
  basicstyle=\ttfamily\small,
  keywordstyle=\color{blue!70!black},
  commentstyle=\color{green!50!black},
  stringstyle=\color{red!60!black},
  backgroundcolor=\color{gray!5},
  frame=single,
  framerule=0.4pt,
  rulecolor=\color{gray!50},
  breaklines=true,
  showstringspaces=false,
  tabsize=4,
  aboveskip=0.8\baselineskip,
  belowskip=0.8\baselineskip
}

\title{Tool Exploration: \texttt{edgar-crawler} for SEC Filing Extraction}
\author{Marco Zhang}
\date{February 2026}

\begin{document}
\maketitle

% ---------------------------------------------------------------------------
% Section 1: Introduction / Goal
% ---------------------------------------------------------------------------
\section{Introduction}

Our goal: download SEC filings in bulk from the authoritative source (SEC EDGAR)
and extract structured, item-level text---not raw HTML blobs---for use across
multiple research projects. Three requirements shape our tool choice.

\subsection*{Batch Downloading}

SEC EDGAR hosts millions of filings. Downloading them one at a time is impractical:

\begin{lstlisting}
# Manual, one at a time
download("AAPL", "10-K", "2023")
download("AAPL", "10-K", "2022")
download("AAPL", "10-K", "2021")
# ... repeat 1000 times for all companies/years
\end{lstlisting}

We need a tool that automates bulk download:

\begin{lstlisting}
# Automated bulk download
download_all(
    years=[1996, 2024],
    forms=["10-K", "10-Q"],
    companies="ALL"  # or specific list
)
# Downloads 100,000+ filings automatically
\end{lstlisting}

\subsection*{Multi-Filter Support}

Some tools (e.g., \texttt{py-edgar}) allow filtering by only one dimension at a
time---CIK or form type, but not both. We need to combine filters: specific
companies + date range + form type + quarter:

\begin{lstlisting}
# Single filter (py-edgar):
download(cik="0000320193")  # Only Apple
# But can't combine: "Apple + Microsoft, 2020-2023, 10-Ks only"

# Multiple filters (edgar-crawler):
download(
    tickers=["AAPL", "MSFT", "GOOGL"],
    years=[2020, 2023],
    forms=["10-K"],
    quarters=[4]
)
\end{lstlisting}

\subsection*{Structured Extraction}

Most tools return raw HTML---one giant blob per filing. For research, we need
item-level text (e.g., Item~1A Risk Factors, Item~7 MD\&A) already separated:

\begin{lstlisting}[language={}]
# Without section extraction: one giant blob
<!DOCTYPE html><html><head>...
ITEM 1. BUSINESS...50 pages...
ITEM 1A. RISK FACTORS...80 pages...
[Embedded tables, HTML tags, images, etc.]
\end{lstlisting}

\begin{lstlisting}[language={}]
# With section extraction: structured JSON
{"item_1": "...", "item_1A": "...", "item_7": "..."}
\end{lstlisting}

Only two tools satisfy all three requirements: \textbf{edgar-crawler}
(open-source, Python) and \textbf{sec-api.io} (paid web service). Since
\texttt{sec-api.io} is a paid service, we explore \texttt{edgar-crawler}
specifically below.

% ---------------------------------------------------------------------------
% Section 2: How We Use It
% ---------------------------------------------------------------------------
\section{How We Use It}

Each project simply uses its own \texttt{config.json} pointing to its own
folders. The tool is stateless---all state lives in the config and output
directories:

\begin{lstlisting}[language={}]
project-A/config.json -> project-A/RAW_FILINGS/ -> project-A/EXTRACTED_FILINGS/
project-B/config.json -> project-B/RAW_FILINGS/ -> project-B/EXTRACTED_FILINGS/
\end{lstlisting}

No code changes needed between projects---just swap the config. Each
\texttt{config.json} specifies two parts corresponding to the two-step
workflow:

\subsection*{Step 1: Download Raw Filings (\texttt{download\_filings})}

Specify which filings to fetch---company tickers, year range, form types,
and quarters---then run \texttt{python download\_filings.py}. This downloads
raw HTML/txt files from EDGAR and produces a metadata CSV.

\begin{lstlisting}[language={}]
"download_filings": {
    "start_year": 2023,
    "end_year": 2024,
    "quarters": [1, 2, 3, 4],
    "filing_types": ["10-K", "10-Q"],
    "cik_tickers": ["AAPL"],          // or [] for all companies
    "user_agent": "Your Name (your@email.edu)",
    "raw_filings_folder": "RAW_FILINGS",
    "indices_folder": "INDICES",
    "filings_metadata_file": "FILINGS_METADATA.csv",
    "skip_present_indices": true
}
\end{lstlisting}

\subsection*{Step 2: Extract Items (\texttt{extract\_items})}

Parse each raw filing into a structured JSON with one key per item. Run
\texttt{python extract\_items.py}.

\begin{lstlisting}[language={}]
"extract_items": {
    "raw_filings_folder": "RAW_FILINGS",
    "extracted_filings_folder": "EXTRACTED_FILINGS",
    "filings_metadata_file": "FILINGS_METADATA.csv",
    "filing_types": ["10-K", "10-Q"],
    "items_to_extract": [],            // [] = all items
    "remove_tables": true,
    "skip_extracted_filings": true
}
\end{lstlisting}

Output: \texttt{EXTRACTED\_FILINGS/} contains one JSON per filing. For a 10-K:

\begin{lstlisting}[language={}]
{"item_1": "Apple Inc. designs...",
 "item_1A": "The following discussion of risk factors...",
 "item_7": "This MD&A section...", ...}
\end{lstlisting}

For a 10-Q:

\begin{lstlisting}[language={}]
{"part_1": "(full Part 1 blob)",
 "part_1_item_2": "Management's Discussion and Analysis...",
 "part_2": "(full Part 2 blob)",
 "part_2_item_1": "Legal Proceedings...", ...}
\end{lstlisting}

% ---------------------------------------------------------------------------
% Section 3: Configuration & Limitations
% ---------------------------------------------------------------------------
\section{Configuration \& Limitations}

The tool's authors provide extensive checks for 10-K, 10-Q, and 8-K extraction,
which produce good structured results. Everything sounds perfect---and overall,
it is reliable. But some limitations emerge from how the configuration works. To
understand them, we look at the two config stages: file-level (Step~1) and
item-level (Step~2). We skip all other config options; see the full reference at
\url{https://github.com/lefterisloukas/edgar-crawler}.

% --- 3.1 File-Level ---
\subsection{File-Level Extraction (\texttt{filing\_types})}

The \texttt{filing\_types} list controls which form types are downloaded. The
default is \texttt{["10-K", "8-K", "10-Q"]}. Importantly, these three are the
only types the authors have verified with multiple checks---the GitHub README
states: ``We currently support the structured JSON output for 10-K, 10-Q and 8-K
documents.'' This raises a few things to be aware of.

EDGAR uses exact string matching on form types
(\texttt{df.Type.isin(filing\_types)}). Specifying just \texttt{"10-K"} misses
all variants. We fetched the real EDGAR index for Q1+Q2 2005 (589,394 total
filings) and confirmed: \texttt{.isin(["10-K"])} captures only 7,726 out of
14,643 annual-report-related filings---\textbf{missing 47\%}. Nine distinct
variants exist:

\begin{table}[h]
\centering
\begin{tabular}{lrl}
\toprule
Type & Count & Description \\
\midrule
10-K     & 7,726 & Standard annual report \\
10KSB    & 2,561 & Small business annual (discontinued 2009) \\
NT 10-K  & 2,040 & Late filing notification \\
10-K/A   & 1,540 & Amended annual report \\
10KSB/A  & 687   & Amended small business annual \\
NTN 10K  & 49    & Non-timely notification \\
NT 10-K/A & 31   & Late filing notification for amendment \\
10-KT    & 7     & Transition period annual \\
10-KT/A  & 2     & Amended transition period \\
\bottomrule
\end{tabular}
\caption{Annual-report-related filing types in EDGAR Q1+Q2 2005 index.}
\end{table}

All 9 variants are fetchable (HTTP 200 confirmed). Must explicitly list all
desired variants in the config. If we want consistent extraction from any of
these, we need to further understand their document structure and potentially
update the extraction code.

% --- 3.2 Item-Level ---
\subsection{Item-Level Extraction (\texttt{extract\_items.py})}

The tool maintains a static list of all items to extract. For 10-K, this is
\texttt{item\_list\_10k}:

\begin{lstlisting}[language={}]
["1", "1A", "1B", "1C", "2", "3", "4", "5", "6", "7", "7A", "8",
 "9", "9A", "9B", "9C", "10", "11", "12", "13", "14", "15", "16"]
\end{lstlisting}

% --- 3.2.1 Stale Item Config ---
\subsubsection{Stale Item Config}

This list is a superset of all known 10-K items across all eras---but it
requires manual updates whenever the SEC introduces, eliminates, or repurposes
an item. Two recent examples: Item~1C (Cybersecurity) was mandated for fiscal
years ending after Dec~15, 2023; Item~6 (Selected Financial Data) was eliminated
in 2021, though companies still include it as ``[RESERVED]''.

We ran the tool's own test suite and all 3 tests failed. The primary cause: the
author added Item~1C to \texttt{item\_list\_10k} but didn't regenerate the test
fixtures, so every expected JSON was missing the \texttt{item\_1C} key. We
raised a PR on the original repo to fix the stale fixtures. The deeper lesson:
the item set requires active management as SEC rules evolve.

To understand what to expect from the extraction, we built an item evolution
table from the 62 10-K test fixtures (1994--2023):

\begin{table}[h]
\centering
\begin{tabular}{lll}
\toprule
Item & First Appears & Notes \\
\midrule
1--6, 8--14          & 1994 (earliest)  & Consistently present \\
7A                   & 1997             & Quantitative risk disclosures \\
1A (Risk Factors)    & $\sim$2004       & SOX-era; empty pre-2004 \\
9A (Controls)        & $\sim$2004       & SOX-era; empty pre-2004 \\
9B (Other Info)      & $\sim$2004--2006 & Appears with SOX items \\
1B (Staff Comments)  & $\sim$2006       & Spotty early on \\
15 (Exhibits)        & 1994             & Consistent \\
16 (10-K Summary)    & Rarely           & Empty in nearly all fixtures \\
1C (Cybersecurity)   & 2023+            & SEC mandate FY ending Dec 15, 2023+ \\
9C (Foreign Jurisd.) & $\sim$2020+      & HFCAA \\
6 (Consolidated Data)& 1994--2023       & Still present despite SEC elimination 2021 \\
\bottomrule
\end{tabular}
\caption{10-K item evolution across 62 test fixtures (1994--2023).}
\end{table}

Three eras emerge. \textbf{Pre-SOX (1994--2003)}---no 1A/9A/9B; risk factor
research is impossible at item level. \textbf{Post-SOX (2004--2022)}---full set
except 1C/9C/16; the ``standard'' era for most research. \textbf{Modern
(2023+)}---1C added, 6 eliminated; structural breaks for longitudinal studies.
Key takeaway: empty items pre-2004 are mostly regulatory non-existence, not
extraction failure.

% --- 3.2.2 Combined Items Detection ---
\subsubsection{Combined Items Detection}

Companies sometimes combine items in their filing (e.g., ``Items 9, 9A, and
9B''). When this happens, all content is assigned to the first item only; the
subsequent items appear empty. The tool does not detect or flag this. Needs
post-processing detection for affected items.

% --- 3.2.3 10-Q Time-Series Consistency ---
\subsubsection{10-Q Time-Series Consistency (Part-Level vs.\ Item-Level)}

The tool extracts 10-Q at two levels: the whole part blob (\texttt{part\_1},
\texttt{part\_2}) and individual items within each part
(\texttt{part\_1\_item\_1}, \texttt{part\_1\_item\_2}, etc.). For 6 out of 200
(3\%) test fixtures---all from 1993--1995---the tool successfully grabs the
entire Part~1 text (6K--14K chars), but when it tries to split into individual
items, it gets nothing: \texttt{part\_1\_item\_1} through
\texttt{part\_1\_item\_4} are all empty. Part~2 items extract correctly in all
cases.

The cause: old \texttt{.txt} filings (pre-$\sim$1996) don't have explicit
``ITEM~1'', ``ITEM~2'' headers within Part~I---they go straight from ``PART~I -
FINANCIAL INFORMATION'' into the content. The regex finds nothing. In the
real-world EDGAR population (not curated test fixtures), the failure rate is
likely higher than 3\%. If you need item-level 10-Q data (e.g., MD\&A = Part~1
Item~2) from pre-$\sim$2000 filings, fall back to the \texttt{part\_1} blob and
parse it yourself.

% --- 3.3 Spot Check ---
\subsection{Spot Check}

We manually verified 5 filings across eras from test fixture expected JSONs,
plus one modern 10-Q:

\begin{table}[h]
\centering
\small
\begin{tabular}{lllll}
\toprule
Filing & Era & Format & Extracted & Notes \\
\midrule
Turner Broadcasting 1993 10-K & Pre-2000 & .txt & 14/22 &
  1A/7A/9A empty (pre-SOX). \\
Internet America 2003 10-K & Early HTML & .htm & 17/22 &
  Has 9A. Missing 1A/1B. \\
Walt Disney 2012 10-K & Post-SOX & .htm & 19/22 &
  1A = 28K, 7 = 75K chars. \\
FedEx 2023 10-K & Modern & .htm & 22/22 &
  All items. Item 6 = ``[RESERVED]''. \\
Zurn Industries 1993 10-Q & Pre-2000 & .txt & Part 1: \textbf{failed} &
  Blob = 11K but item keys empty. \\
Microsoft 2024 10-Q & Modern & .htm & Full success &
  Part 1: 4 items, Part 2: 5 items. \\
\bottomrule
\end{tabular}
\caption{Spot-check results across filing eras.}
\end{table}

Interpretation: (1)~10-K extraction is reliable across all eras.
(2)~10-Q has a pre-2000 Part~1 failure mode. (3)~Modern extraction quality is
excellent.

% --- Summary Table ---
\subsection*{Summary}

\begin{table}[h]
\centering
\begin{tabular}{clll}
\toprule
\# & Finding & Impact & Ref. \\
\midrule
1 & 10-Q Part~1 item extraction fails on old .txt filings & High & \S3.2.3 \\
2 & \texttt{filing\_types} exact matching excludes 47\% of variants & High & \S3.1 \\
3 & Item structural breaks across three decades & High & \S3.2.1 \\
4 & Item set requires manual tracking as SEC changes rules & Med-High & \S3.2.1 \\
5 & Combined items not detected & Medium & \S3.2.2 \\
6 & 10-K extraction reliable across all eras (1993--2024) & Positive & \S3.3 \\
7 & Modern extraction (2010+) quality is excellent & Positive & \S3.3 \\
\bottomrule
\end{tabular}
\caption{Summary of findings.}
\end{table}

% ---------------------------------------------------------------------------
% Section 4: Other Minor Things & Next Steps
% ---------------------------------------------------------------------------
\section{Other Minor Things \& Next Steps}

\subsection{Minor Things}

\begin{itemize}
  \item \textbf{Table removal impact on MD\&A (Item~7).}
    Tables embedded in narrative sections like MD\&A are removed by default
    (\texttt{remove\_tables: true}). This is fine for pure text analysis but can
    be toggled off if tables are needed.

  \item \textbf{10-Q part separation fragility.}
    The two-part separation uses 4 heuristic bug-detection cases. Failure modes
    include ToC contamination and arbitrary thresholds. Known limitation; can't
    do better for now.

  \item \textbf{Fiscal year alignment.}
    The \texttt{quarters} config filters by which calendar quarter the filing was
    \emph{submitted} to EDGAR, not the period it reports on. Companies have
    different fiscal year-ends (e.g., Apple = September, most others = December),
    so the same calendar quarter contains filings for different reporting periods.
    Always filter by \texttt{period\_of\_report} in the metadata for
    cross-sectional comparisons.

  \item \textbf{Survivorship bias.}
    Could happen if downloading only a subset of companies via ticker lists
    (tickers of delisted firms may not resolve). Avoided by using CIK identifiers
    or downloading the full universe (\texttt{cik\_tickers: []}).
\end{itemize}

\subsection{Next Steps}

If we decide to go with this tool:
\begin{enumerate}
  \item Set up a multi-project environment with clear hierarchy and markdown
    instruction files that feed Claude Code for automated use (preliminary
    hierarchy already drafted). Since each project uses its own
    \texttt{config.json} (Section~2), different projects can specify different
    filing types, item sets, and time ranges as needed---e.g., one project
    downloads only 10-K, another only 10-Q.
  \item Encode the findings from this report as project-level instructions for
    the Claude agent.
\end{enumerate}

\end{document}
